% GENERATED FILE. Edit canonical lessons, then run npm run generate:books.
% canonical-source-hash: fnv1a64:11572c6ed590ae83
% canonical-lessons: AR-C322-qala, AR-C322-hisn, AR-C322-sur, AR-C322-hasara, AR-C322-dafaa-defend

\chapter{Castles and walls}
\label{ch:ar-castles-and-walls}
\hlchaptermodality{322}

\begin{chapteropening}
\textbf{By the end of this chapter:} \emph{I can say a castle, a fort, a city wall, to besiege and to defend.}

Complete the last lesson of chapter 322: I can say a castle, a fort, a city wall, to besiege and to defend.
\end{chapteropening}

\section[\texorpdfstring{\ar{قلعة} (qalʿa)}{قلعة (qalʿa)}]{\ar{قلعة} (qalʿa) — a castle, a citadel}

\label{lesson:AR-C322-qala}



\begin{quote}
Before the new one: say the Arabic for to rob, then the Arabic for to loot.
\end{quote}

\subsection*{You'll want to know: \ar{قلعة}}
\textbf{\ar{قلعة}} — \emph{qalʿa} — \textquotedblleft{}a castle, a citadel\textquotedblright{}.

\begin{quote}
\ar{صعدنا} \ar{إلى} \ar{قلعة} \ar{على} \ar{الجبل}.
\end{quote}

\begin{quote}
\emph{ṣaʿidnā ilā qalʿa ʿalā al-jabal.}
\end{quote}

\begin{quote}
We climbed up to a castle on the mountain.
\end{quote}

Plural: \ar{قلاع}, \emph{qilāʿ}.

\begin{cousinweb}
\textbf{\ar{قلعة}} (\emph{qalʿa}) is usually filed under the native Arabic root \emph{q-l-ʿ}, \textquotedblleft{}to pull out by the roots\textquotedblright{} (Lane). Lane's article does not give the castle sense, and how a fortress took the name is not settled.
\end{cousinweb}

\begin{grammarlens}[title={What you've built}]
One more word for telling what happened.
\end{grammarlens}

\subsection*{Your turn}
\begin{itemize}
\raggedright
  \item \emph{Say it:} \emph{qalʿa}
  \item \emph{Say it:} \emph{qalʿa}, once more
  \item \emph{Say it:} \emph{ṣaʿidnā ilā qalʿa ʿalā al-jabal.} — We climbed up to a castle on the mountain.
  \item \emph{Recall:} say the Arabic for to rob, then the Arabic for to loot, then say \emph{qalʿa} again
\end{itemize}

\begin{tcolorbox}[breakable,colback=teal!4,colframe=teal!35!black,title={Before you move on}]
What does \ar{قلعة} mean? (\textquotedblleft{}A castle, a citadel\textquotedblright{}.) Say it once more.
\end{tcolorbox}
\section[\texorpdfstring{\ar{حصن} (ḥiṣn)}{حصن (ḥiṣn)}]{\ar{حصن} (ḥiṣn) — a fort, a stronghold}

\label{lesson:AR-C322-hisn}



\begin{quote}
Before the new one: say the Arabic for to loot, then the Arabic for a castle.
\end{quote}

\subsection*{You'll want to know: \ar{حصن}}
\textbf{\ar{حصن}} — \emph{ḥiṣn} — \textquotedblleft{}a fort, a stronghold\textquotedblright{}.

\begin{quote}
\ar{بنى} \ar{الجنود} \ar{حصنا} \ar{قرب} \ar{البحر}.
\end{quote}

\begin{quote}
\emph{banā al-junūd ḥiṣnan qurba al-baḥr.}
\end{quote}

\begin{quote}
The soldiers built a fort near the sea.
\end{quote}

Plural: \ar{حصون}, \emph{ḥuṣūn}.

\begin{cousinweb}
\textbf{\ar{حصن}} (\emph{ḥiṣn}) is built on the native Arabic root \emph{ḥ-ṣ-n} (Lane), \textquotedblleft{}to be inaccessible, well guarded\textquotedblright{}; the taught \emph{ḥiṣān}, a horse, is the same root.
\end{cousinweb}

\begin{grammarlens}[title={What you've built}]
One more word for telling what happened.
\end{grammarlens}

\subsection*{Your turn}
\begin{itemize}
\raggedright
  \item \emph{Say it:} \emph{ḥiṣn}
  \item \emph{Say it:} \emph{ḥiṣn}, once more
  \item \emph{Say it:} \emph{banā al-junūd ḥiṣnan qurba al-baḥr.} — The soldiers built a fort near the sea.
  \item \emph{Recall:} say the Arabic for to loot, then the Arabic for a castle, then say \emph{ḥiṣn} again
\end{itemize}

\begin{tcolorbox}[breakable,colback=teal!4,colframe=teal!35!black,title={Before you move on}]
What does \ar{حصن} mean? (\textquotedblleft{}A fort, a stronghold\textquotedblright{}.) Say it once more.
\end{tcolorbox}
\section[\texorpdfstring{\ar{سور} (sūr)}{سور (sūr)}]{\ar{سور} (sūr) — a city wall}

\label{lesson:AR-C322-sur}



\begin{quote}
Before the new one: say the Arabic for a castle, then the Arabic for a fort.
\end{quote}

\subsection*{You'll want to know: \ar{سور}}
\textbf{\ar{سور}} — \emph{sūr} — \textquotedblleft{}a city wall\textquotedblright{}.

\begin{quote}
\ar{حول} \ar{المدينة} \ar{القديمة} \ar{سور} \ar{طويل}.
\end{quote}

\begin{quote}
\emph{ḥawla al-madīna al-qadīma sūr ṭawīl.}
\end{quote}

\begin{quote}
There is a long wall round the old city.
\end{quote}

Plural: \ar{أسوار}, \emph{aswār}.

A wall round a garden or a school is also \ar{سور}; the wall of a room is \ar{جدار}.

\begin{cousinweb}
\textbf{\ar{سور}} (\emph{sūr}) is built on the native Arabic root \emph{s-w-r} (Lane), where Lane gives \emph{sūr} as the wall of a town, and \emph{sāra}, to climb over it.
\end{cousinweb}

\begin{grammarlens}[title={What you've built}]
One more word for telling what happened.
\end{grammarlens}

\subsection*{Your turn}
\begin{itemize}
\raggedright
  \item \emph{Say it:} \emph{sūr}
  \item \emph{Say it:} \emph{sūr}, once more
  \item \emph{Say it:} \emph{ḥawla al-madīna al-qadīma sūr ṭawīl.} — There is a long wall round the old city.
  \item \emph{Recall:} say the Arabic for a castle, then the Arabic for a fort, then say \emph{sūr} again
\end{itemize}

\begin{tcolorbox}[breakable,colback=teal!4,colframe=teal!35!black,title={Before you move on}]
What does \ar{سور} mean? (\textquotedblleft{}A city wall\textquotedblright{}.) Say it once more.
\end{tcolorbox}
\section[\texorpdfstring{\ar{حاصر} (ḥāṣara)}{حاصر (ḥāṣara)}]{\ar{حاصر} (ḥāṣara) — to besiege, to surround}

\label{lesson:AR-C322-hasara}



\begin{quote}
Before the new one: say the Arabic for a fort, then the Arabic for a city wall.
\end{quote}

\subsection*{You'll want to know: \ar{حاصر}}
\textbf{\ar{حاصر}} — \emph{ḥāṣara} — \textquotedblleft{}to besiege, to surround\textquotedblright{}.

\begin{quote}
\ar{حاصر} \ar{الجيش} \ar{القلعة} \ar{شهرا}.
\end{quote}

\begin{quote}
\emph{ḥāṣara al-jaysh al-qalʿa shahran.}
\end{quote}

\begin{quote}
The army besieged the castle for a month.
\end{quote}

Present: \ar{يحاصر}, \emph{yuḥāṣir}.

\begin{cousinweb}
\textbf{\ar{حاصر}} (\emph{ḥāṣara}) is built on the native Arabic root \emph{ḥ-ṣ-r} (Lane), \textquotedblleft{}to confine, to hem in\textquotedblright{}. It is Form III, \textquotedblleft{}to hem in an enemy\textquotedblright{}.
\end{cousinweb}

\begin{grammarlens}[title={What you've built}]
One more word for telling what happened.
\end{grammarlens}

\subsection*{Your turn}
\begin{itemize}
\raggedright
  \item \emph{Say it:} \emph{ḥāṣara}
  \item \emph{Say it:} \emph{ḥāṣara}, once more
  \item \emph{Say it:} \emph{ḥāṣara al-jaysh al-qalʿa shahran.} — The army besieged the castle for a month.
  \item \emph{Recall:} say the Arabic for a fort, then the Arabic for a city wall, then say \emph{ḥāṣara} again
\end{itemize}

\begin{tcolorbox}[breakable,colback=teal!4,colframe=teal!35!black,title={Before you move on}]
What does \ar{حاصر} mean? (\textquotedblleft{}To besiege, to surround\textquotedblright{}.) Say it once more.
\end{tcolorbox}
\section[\texorpdfstring{\ar{دافع} (dāfaʿa)}{دافع (dāfaʿa)}]{\ar{دافع} (dāfaʿa) — to defend}

\label{lesson:AR-C322-dafaa-defend}



\begin{quote}
Before the new one: say the Arabic for a city wall, then the Arabic for to besiege.
\end{quote}

\subsection*{You'll want to know: \ar{دافع}}
\textbf{\ar{دافع}} — \emph{dāfaʿa} — \textquotedblleft{}to defend\textquotedblright{}.

\begin{quote}
\ar{دافع} \ar{الجنود} \ar{عن} \ar{المدينة}.
\end{quote}

\begin{quote}
\emph{dāfaʿa al-junūd ʿan al-madīna.}
\end{quote}

\begin{quote}
The soldiers defended the city.
\end{quote}

Present: \ar{يدافع}, \emph{yudāfiʿ}.

It takes \ar{عن}: \ar{دافع} \ar{عن} \ar{صديقه}.

\begin{cousinweb}
\textbf{\ar{دافع}} (\emph{dāfaʿa}) is built on the native Arabic root \emph{d-f-ʿ} (Lane), \textquotedblleft{}to push\textquotedblright{}, the root of the taught \emph{dafaʿa}. It is Form III, \textquotedblleft{}to push back against someone\textquotedblright{}: to defend.
\end{cousinweb}

\begin{grammarlens}[title={What you've built}]
One more word for telling what happened.
\end{grammarlens}

\subsection*{Your turn}
\begin{itemize}
\raggedright
  \item \emph{Say it:} \emph{dāfaʿa}
  \item \emph{Say it:} \emph{dāfaʿa}, once more
  \item \emph{Say it:} \emph{dāfaʿa al-junūd ʿan al-madīna.} — The soldiers defended the city.
  \item \emph{Recall:} say the Arabic for a city wall, then the Arabic for to besiege, then say \emph{dāfaʿa} again
\end{itemize}

\begin{tcolorbox}[breakable,colback=teal!4,colframe=teal!35!black,title={Before you move on}]
What does \ar{دافع} mean? (\textquotedblleft{}To defend\textquotedblright{}.) Say it once more.
\end{tcolorbox}
